\section{Conclusiones}

\subsection{Experimento 1: efecto del largo del texto}

Al aumentar el texto de 1 a 10 MB, los tiempos de fuerza bruta, KMP y
Rabin--Karp crecieron entre 7 y 8 veces, comportamiento compatible con un
crecimiento aproximadamente lineal en el rango estudiado. Fuerza bruta
presentó los menores tiempos entre los métodos sin índice sobre estas
entradas, lo que muestra que una mejor cota de peor caso no garantiza un
menor tiempo en cada conjunto de datos. La proporcionalidad no fue exacta;
los resultados no permiten atribuir esas desviaciones a una causa concreta.

Las consultas de Suffix Array se mantuvieron por debajo de un microsegundo,
sin crecimiento sostenido. Esto no demuestra tiempo constante ni que un
texto mayor favorezca la búsqueda: su cota es \(O(m\log n)\), y multiplicar
el texto por diez cambia poco su logaritmo. Para 10 MB, la consulta promedió
\(0{,}532\,\mu\mathrm{s}\), pero construir el índice requirió \(34{,}57\) s.
El índice persistente ocupó 90 MB y el pico total de construcción alcanzó
410 MB; ambos crecieron linealmente con el texto. El pico ya incluye las
reservas persistentes presentes en ese instante, por lo que no se suman
ambas magnitudes. La ventaja de consulta tiene, por tanto, un costo inicial
y espacial que debe evaluarse por separado.

\subsection{Experimento 2: efecto del largo del patrón}

Al extender el patrón de 8 a 256 bytes, no se observó un aumento sostenido
del tiempo. Para KMP, esto es coherente con \(O(n+m)\): frente a un texto
fijo de diez millones de bytes, la preparación de patrones de hasta
256 bytes tiene un peso reducido. Rabin--Karp también recorre prácticamente
la misma cantidad de ventanas. En fuerza bruta, muchas comparaciones pueden
terminar antes de examinar el patrón completo, por lo que su cota
\(O(nm)\) no exige un crecimiento proporcional a \(m\) en estas entradas.

La caída observada para fuerza bruta con patrones de 32 bytes no permite
concluir que esa longitud sea óptima: entre longitudes cambian también los
patrones y sus prefijos. Suffix Array tampoco mostró una penalización
sostenida al aumentar \(m\), sin que ello invalide su cota teórica. En
memoria, KMP sí incrementa la tabla auxiliar con el patrón, mientras el
índice persistente permanece fijo porque el texto no cambia. La conclusión
se limita, así, a una baja sensibilidad temporal en el rango evaluado.

\subsection{Experimento 3: efecto del número de ocurrencias}

En el corpus natural no apareció una relación proporcional clara entre
frecuencia y tiempo. Los métodos sin índice recorren el texto incluso si
el patrón está ausente; además, patrones con distintas frecuencias también
difieren en contenido. Por ello, no puede atribuirse toda variación a
\(z\). Suffix Array localiza los límites del intervalo de coincidencias y
los resta, sin visitar cada resultado: contar no exige el mismo trabajo
que devolver todas sus posiciones.

Los casos \(T=A^n\), con \(P=A\) y \(P=B\), contrastan presencia máxima y
ausencia, pero no establecen que más coincidencias siempre implique mayor
tiempo. Por ejemplo, para 10 MB KMP promedió aproximadamente 19 ms con
\(P=A\) y 86 ms con \(P=B\). Al ser \(m=1\), los tres métodos sin índice
mantienen costo lineal en \(n\) en ambos casos; estas pruebas no reproducen
el peor caso cuadrático de fuerza bruta con patrones de longitud creciente.

\clearpage
\subsection{Conclusión general}

Los experimentos muestran que la selección de un método no debe basarse en
una clasificación única de rapidez. Es necesario distinguir el costo de
una consulta aislada del costo total de preparar y reutilizar una
estructura, así como separar las garantías asintóticas del desempeño
observado en un corpus concreto. Una curva casi horizontal en un gráfico
logarítmico no demuestra complejidad constante, y una implementación con
peor cota teórica puede resultar competitiva cuando las entradas no
presentan su caso desfavorable.

Para pocas consultas o textos que cambian frecuentemente, los métodos sin
índice evitan una construcción costosa y la reserva persistente asociada.
Fuerza bruta constituye una alternativa sencilla y fue competitiva sobre
el texto natural evaluado. Sin embargo, no ofrece la garantía
\(O(n+m)\) de KMP, que resulta preferible cuando interesa limitar el costo
ante entradas desfavorables o repetitivas. Su mayor tiempo en parte de
estos experimentos no elimina esa ventaja: el caso promedio observado y
la protección frente al peor caso responden a necesidades distintas.

Para muchas consultas sobre un texto estático, Suffix Array puede
compensar su construcción mediante búsquedas mucho más rápidas, siempre
que haya memoria suficiente. Si \(Q\) es la cantidad de consultas, una
comparación simplificada considera
\[
    T_{\mathrm{SA}}(Q)=T_{\text{construcción}}+Q\,t_{\mathrm{SA}},
    \qquad
    T_{\text{sin índice}}(Q)=Q\,t_{\text{sin índice}}.
\]
Este modelo ilustra la amortización, pero no determina un umbral universal:
los tiempos dependen de los patrones, del equipo y de la posibilidad de
reutilizar el índice. Tampoco representa una medición de una carga mixta
de consultas o de actualizaciones del texto.

Rabin--Karp, con esta implementación y \(q=101\), no mostró una ventaja
temporal frente a las otras alternativas evaluadas. No corresponde
generalizar ese resultado a todas sus variantes ni atribuirlo únicamente
a colisiones, cuyo número no se midió. Cambiar la configuración del hash
y registrar las verificaciones permitiría estudiar mejor ese comportamiento.

El alcance de estas conclusiones está limitado por un solo corpus natural,
los tamaños seleccionados y los casos sintéticos específicos. Las
32 repeticiones describen variabilidad temporal, pero no reemplazan la
diversidad de entradas; sus desviaciones estándar tampoco prueban
significancia estadística. Además, se midió conteo, no recuperación de
posiciones, y memoria dinámica, no memoria total del proceso. Como extensión,
sería útil evaluar otros corpus, patrones repetitivos de longitud creciente
y cargas de muchas consultas que incluyan la construcción del índice.

En síntesis, no existe una alternativa universalmente superior. La decisión
debe equilibrar cantidad de consultas, estabilidad del texto, memoria
disponible y garantías requeridas, utilizando la teoría para anticipar
riesgos y la experimentación para estimar los costos efectivos del
escenario de uso.
